%!TeX root=../persuasiontop.tex
\chapter[Chapter \thechapter]{}

\begin{figure}[t!]
\centering
\includegraphics[width=\linewidth]{nouveau03}
\end{figure}

\lettrine[lraise=0.3,ante=`]{I}{} must take leave to observe, Sir Walter,' said Mr Shepherd one morning at Kellynch Hall, as he laid down the newspaper, <that the present juncture is much in our favour. This peace will be turning all our rich naval officers ashore. They will be all wanting a home. Could not be a better time, Sir Walter, for having a choice of tenants, very responsible tenants. Many a noble fortune has been made during the war. If a rich admiral were to come in our way, Sir Walter—>

<He would be a very lucky man, Shepherd,> replied Sir Walter; <that's all I have to remark. A prize indeed would Kellynch Hall be to him; rather the greatest prize of all, let him have taken ever so many before; hey, Shepherd?>

Mr Shepherd laughed, as he knew he must, at this wit, and then added—

<I presume to observe, Sir Walter, that, in the way of business, gentlemen of the navy are well to deal with. I have had a little knowledge of their methods of doing business; and I am free to confess that they have very liberal notions, and are as likely to make desirable tenants as any set of people one should meet with. Therefore, Sir Walter, what I would take leave to suggest is, that if in consequence of any rumours getting abroad of your intention; which must be contemplated as a possible thing, because we know how difficult it is to keep the actions and designs of one part of the world from the notice and curiosity of the other; consequence has its tax; I, John Shepherd, might conceal any family-matters that I chose, for nobody would think it worth their while to observe me; but Sir Walter Elliot has eyes upon him which it may be very difficult to elude; and therefore, thus much I venture upon, that it will not greatly surprise me if, with all our caution, some rumour of the truth should get abroad; in the supposition of which, as I was going to observe, since applications will unquestionably follow, I should think any from our wealthy naval commanders particularly worth attending to; and beg leave to add, that two hours will bring me over at any time, to save you the trouble of replying.>

Sir Walter only nodded. But soon afterwards, rising and pacing the room, he observed sarcastically—

<There are few among the gentlemen of the navy, I imagine, who would not be surprised to find themselves in a house of this description.>

<They would look around them, no doubt, and bless their good fortune,> said Mrs Clay, for Mrs Clay was present: her father had driven her over, nothing being of so much use to Mrs Clay's health as a drive to Kellynch: <but I quite agree with my father in thinking a sailor might be a very desirable tenant. I have known a good deal of the profession; and besides their liberality, they are so neat and careful in all their ways! These valuable pictures of yours, Sir Walter, if you chose to leave them, would be perfectly safe. Everything in and about the house would be taken such excellent care of! The gardens and shrubberies would be kept in almost as high order as they are now. You need not be afraid, Miss Elliot, of your own sweet flower gardens being neglected.>

<As to all that,> rejoined Sir Walter coolly, <supposing I were induced to let my house, I have by no means made up my mind as to the privileges to be annexed to it. I am not particularly disposed to favour a tenant. The park would be open to him of course, and few navy officers, or men of any other description, can have had such a range; but what restrictions I might impose on the use of the pleasure-grounds, is another thing. I am not fond of the idea of my shrubberies being always approachable; and I should recommend Miss Elliot to be on her guard with respect to her flower garden. I am very little disposed to grant a tenant of Kellynch Hall any extraordinary favour, I assure you, be he sailor or soldier.>

After a short pause, Mr Shepherd presumed to say—

<In all these cases, there are established usages which make everything plain and easy between landlord and tenant. Your interest, Sir Walter, is in pretty safe hands. Depend upon me for taking care that no tenant has more than his just rights. I venture to hint, that Sir Walter Elliot cannot be half so jealous for his own, as John Shepherd will be for him.>

Here Anne spoke—

<The navy, I think, who have done so much for us, have at least an equal claim with any other set of men, for all the comforts and all the privileges which any home can give. Sailors work hard enough for their comforts, we must all allow.>

<Very true, very true. What Miss Anne says, is very true,> was Mr Shepherd's rejoinder, and <Oh! certainly,> was his daughter's; but Sir Walter's remark was, soon afterwards—

<The profession has its utility, but I should be sorry to see any friend of mine belonging to it.>

<Indeed!> was the reply, and with a look of surprise.

<Yes; it is in two points offensive to me; I have two strong grounds of objection to it. First, as being the means of bringing persons of obscure birth into undue distinction, and raising men to honours which their fathers and grandfathers never dreamt of; and secondly, as it cuts up a man's youth and vigour most horribly; a sailor grows old sooner than any other man. I have observed it all my life. A man is in greater danger in the navy of being insulted by the rise of one whose father, his father might have disdained to speak to, and of becoming prematurely an object of disgust himself, than in any other line. One day last spring, in town, I was in company with two men, striking instances of what I am talking of; Lord St Ives, whose father we all know to have been a country curate, without bread to eat; I was to give place to Lord St Ives, and a certain Admiral Baldwin, the most deplorable-looking personage you can imagine; his face the colour of mahogany, rough and rugged to the last degree; all lines and wrinkles, nine grey hairs of a side, and nothing but a dab of powder at top. <In the name of heaven, who is that old fellow?> said I to a friend of mine who was standing near, (Sir Basil Morley). <Old fellow!> cried Sir Basil, <it is Admiral Baldwin. What do you take his age to be?> <Sixty,> said I, <or perhaps sixty-two.> <Forty,> replied Sir Basil, <forty, and no more.> Picture to yourselves my amazement; I shall not easily forget Admiral Baldwin. I never saw quite so wretched an example of what a sea-faring life can do; but to a degree, I know it is the same with them all: they are all knocked about, and exposed to every climate, and every weather, till they are not fit to be seen. It is a pity they are not knocked on the head at once, before they reach Admiral Baldwin's age.>

<Nay, Sir Walter,> cried Mrs Clay, <this is being severe indeed. Have a little mercy on the poor men. We are not all born to be handsome. The sea is no beautifier, certainly; sailors do grow old betimes; I have observed it; they soon lose the look of youth. But then, is not it the same with many other professions, perhaps most other? Soldiers, in active service, are not at all better off: and even in the quieter professions, there is a toil and a labour of the mind, if not of the body, which seldom leaves a man's looks to the natural effect of time. The lawyer plods, quite care-worn; the physician is up at all hours, and travelling in all weather; and even the clergyman—> she stopt a moment to consider what might do for the clergyman;—<and even the clergyman, you know is obliged to go into infected rooms, and expose his health and looks to all the injury of a poisonous atmosphere. In fact, as I have long been convinced, though every profession is necessary and honourable in its turn, it is only the lot of those who are not obliged to follow any, who can live in a regular way, in the country, choosing their own hours, following their own pursuits, and living on their own property, without the torment of trying for more; it is only \textit{their} lot, I say, to hold the blessings of health and a good appearance to the utmost: I know no other set of men but what lose something of their personableness when they cease to be quite young.>

It seemed as if Mr Shepherd, in this anxiety to bespeak Sir Walter's good will towards a naval officer as tenant, had been gifted with foresight; for the very first application for the house was from an Admiral Croft, with whom he shortly afterwards fell into company in attending the quarter sessions at Taunton; and indeed, he had received a hint of the Admiral from a London correspondent. By the report which he hastened over to Kellynch to make, Admiral Croft was a native of Somersetshire, who having acquired a very handsome fortune, was wishing to settle in his own country, and had come down to Taunton in order to look at some advertised places in that immediate neighbourhood, which, however, had not suited him; that accidentally hearing—(it was just as he had foretold, Mr Shepherd observed, Sir Walter's concerns could not be kept a secret,)—accidentally hearing of the possibility of Kellynch Hall being to let, and understanding his (Mr Shepherd's) connection with the owner, he had introduced himself to him in order to make particular inquiries, and had, in the course of a pretty long conference, expressed as strong an inclination for the place as a man who knew it only by description could feel; and given Mr Shepherd, in his explicit account of himself, every proof of his being a most responsible, eligible tenant.

<And who is Admiral Croft?> was Sir Walter's cold suspicious inquiry.

Mr Shepherd answered for his being of a gentleman's family, and mentioned a place; and Anne, after the little pause which followed, added—

<He is a rear admiral of the white. He was in the Trafalgar action, and has been in the East Indies since; he was stationed there, I believe, several years.>

<Then I take it for granted,> observed Sir Walter, <that his face is about as orange as the cuffs and capes of my livery.>

Mr Shepherd hastened to assure him, that Admiral Croft was a very hale, hearty, well-looking man, a little weather-beaten, to be sure, but not much, and quite the gentleman in all his notions and behaviour; not likely to make the smallest difficulty about terms, only wanted a comfortable home, and to get into it as soon as possible; knew he must pay for his convenience; knew what rent a ready-furnished house of that consequence might fetch; should not have been surprised if Sir Walter had asked more; had inquired about the manor; would be glad of the deputation, certainly, but made no great point of it; said he sometimes took out a gun, but never killed; quite the gentleman.

Mr Shepherd was eloquent on the subject; pointing out all the circumstances of the Admiral's family, which made him peculiarly desirable as a tenant. He was a married man, and without children; the very state to be wished for. A house was never taken good care of, Mr Shepherd observed, without a lady: he did not know, whether furniture might not be in danger of suffering as much where there was no lady, as where there were many children. A lady, without a family, was the very best preserver of furniture in the world. He had seen Mrs Croft, too; she was at Taunton with the admiral, and had been present almost all the time they were talking the matter over.

<And a very well-spoken, genteel, shrewd lady, she seemed to be,> continued he; <asked more questions about the house, and terms, and taxes, than the Admiral himself, and seemed more conversant with business; and moreover, Sir Walter, I found she was not quite unconnected in this country, any more than her husband; that is to say, she is sister to a gentleman who did live amongst us once; she told me so herself: sister to the gentleman who lived a few years back at Monkford. Bless me! what was his name? At this moment I cannot recollect his name, though I have heard it so lately. Penelope, my dear, can you help me to the name of the gentleman who lived at Monkford: Mrs Croft's brother?>

But Mrs Clay was talking so eagerly with Miss Elliot, that she did not hear the appeal.

<I have no conception whom you can mean, Shepherd; I remember no gentleman resident at Monkford since the time of old Governor Trent.>

<Bless me! how very odd! I shall forget my own name soon, I suppose. A name that I am so very well acquainted with; knew the gentleman so well by sight; seen him a hundred times; came to consult me once, I remember, about a trespass of one of his neighbours; farmer's man breaking into his orchard; wall torn down; apples stolen; caught in the fact; and afterwards, contrary to my judgement, submitted to an amicable compromise. Very odd indeed!>

After waiting another moment—

<You mean Mr Wentworth, I suppose?> said Anne.

Mr Shepherd was all gratitude.

<Wentworth was the very name! Mr Wentworth was the very man. He had the curacy of Monkford, you know, Sir Walter, some time back, for two or three years. Came there about the year —5, I take it. You remember him, I am sure.>

<Wentworth? Oh! ay, Mr Wentworth, the curate of Monkford. You misled me by the term \textit{gentleman}. I thought you were speaking of some man of property: Mr Wentworth was nobody, I remember; quite unconnected; nothing to do with the Strafford family. One wonders how the names of many of our nobility become so common.>

As Mr Shepherd perceived that this connexion of the Crofts did them no service with Sir Walter, he mentioned it no more; returning, with all his zeal, to dwell on the circumstances more indisputably in their favour; their age, and number, and fortune; the high idea they had formed of Kellynch Hall, and extreme solicitude for the advantage of renting it; making it appear as if they ranked nothing beyond the happiness of being the tenants of Sir Walter Elliot: an extraordinary taste, certainly, could they have been supposed in the secret of Sir Walter's estimate of the dues of a tenant.

It succeeded, however; and though Sir Walter must ever look with an evil eye on anyone intending to inhabit that house, and think them infinitely too well off in being permitted to rent it on the highest terms, he was talked into allowing Mr Shepherd to proceed in the treaty, and authorising him to wait on Admiral Croft, who still remained at Taunton, and fix a day for the house being seen.

Sir Walter was not very wise; but still he had experience enough of the world to feel, that a more unobjectionable tenant, in all essentials, than Admiral Croft bid fair to be, could hardly offer. So far went his understanding; and his vanity supplied a little additional soothing, in the Admiral's situation in life, which was just high enough, and not too high. <I have let my house to Admiral Croft,> would sound extremely well; very much better than to any mere \textit{Mr}\doubleemdash ; a \textit{Mr} (save, perhaps, some half dozen in the nation,) always needs a note of explanation. An admiral speaks his own consequence, and, at the same time, can never make a baronet look small. In all their dealings and intercourse, Sir Walter Elliot must ever have the precedence.

Nothing could be done without a reference to Elizabeth: but her inclination was growing so strong for a removal, that she was happy to have it fixed and expedited by a tenant at hand; and not a word to suspend decision was uttered by her.

Mr Shepherd was completely empowered to act; and no sooner had such an end been reached, than Anne, who had been a most attentive listener to the whole, left the room, to seek the comfort of cool air for her flushed cheeks; and as she walked along a favourite grove, said, with a gentle sigh, <A few months more, and \textit{he}, perhaps, may be walking here.>